% \newpage

\clearpage


\definecolor{mybackcolor}{HTML}{c2eaf7}
\definecolor{myframecolor}{HTML}{8dd4f0}

\section{Appendix}

\subsection{Limitations}
Despite providing a focused and rigorous evaluation of scientific search and reasoning capabilities, \textsc{SciExplore} has several limitations that point toward important directions for future work. 

\noindent\textbf{First, the scale of \textsc{SciExplore} reflects an intentional quality–coverage trade-off.} 
All tasks are manually constructed by domain experts and subjected to strict multi-stage filtering and validation, during which a large number of low-quality or weakly constrained candidates are discarded. 
As a result, the final benchmark consists of just over one hundred high-quality tasks.
While this expert-driven curation ensures strong difficulty control and answer uniqueness, it may limit coverage of the full breadth of scientific disciplines and research styles.
Future work will explore expanding the benchmark while preserving its core quality guarantees.

\noindent\textbf{Second, the current benchmark primarily emphasizes text-centric scientific information sources.}
\textsc{SciExplore} focuses on reasoning over structured databases and scientific literature, which constitute a central component of many research workflows.
However, real-world scientific inquiry often involves additional modalities such as figures, tables, experimental datasets, or code artifacts.
Extending the benchmark to incorporate multi-modal scientific evidence and heterogeneous data formats remains an important direction for future research.

\noindent\textbf{Finally, the maintenance and factual verification of \textsc{SciExplore} currently rely on expert-driven manual updates.}
To ensure correctness and answer uniqueness, all tasks are curated and fact-checked by domain experts against the underlying scientific databases and literature at the time of construction.
While this process guarantees high factual accuracy, it is time-consuming and labor-intensive, making frequent or real-time updates difficult to sustain.
An important direction for future work is to develop automated or semi-automated updating pipelines that can continuously verify, refresh, and expand benchmark content as scientific databases and literature evolve.



% First, the current benchmark contains a relatively limited number of queries (103 expert-curated tasks), which, while sufficient for in-depth qualitative analysis and controlled difficulty design, may not fully capture the long-tail diversity of real-world scientific information-seeking behaviors across domains and subfields. In future iterations, we plan to substantially expand the query set in both scale and topical coverage, incorporating a broader range of scientific disciplines, problem formulations, and difficulty levels to improve statistical robustness and ecological validity. 
% Second, the present evaluation focuses primarily on final task success and structured output correctness; future work could introduce finer-grained process-level metrics to better characterize intermediate search decisions, hypothesis refinement, and evidence verification behaviors. 
% Finally, as scientific tools and databases continue to evolve, we envision extending \textsc{SciExplore} into a continuously updated benchmark that tracks progress in long-horizon scientific agents under realistic, dynamically changing research environments.

\subsection{Models and API Identifiers}
\label{app:models_api_ids}

The models used in our benchmark and their corresponding API identifiers (which may include version or
date information) are listed in Table~\ref{tab:models_api_ids}. The ``Model/Product'' column denotes the
shorthand name used throughout the paper.

It is worth noting that \textit{GPT-5.1} and \textit{GPT-5.1 w/search} share the same API identifier but are
evaluated under different tool configurations (with search disabled and enabled, respectively).
Similarly, \textit{Gemini-3-pro} and \textit{Gemini-3-pro w/search} use the same API identifier with different
search settings.

\begin{table}[t]
\centering
\footnotesize
\resizebox{\columnwidth}{!}{%
\begin{tabular}{ll}
\toprule
\textbf{Model / Product} & \textbf{API Identifier} \\
\midrule
GPT-5.1 & gpt-5.1-2025-11-13 \\
Gemini-2.5-Pro & gemini-2.5-pro \\
Gemini-3-Pro & gemini-3-pro-preview \\
GPT-5.1 w/ search & gpt-5.1-2025-11-13 \\
Gemini3-Pro w/ search & gemini-3-pro-preview \\
Gemini Deep Research & deep-research-pro-preview-12-25 \\
OpenAI Deep Research & o3-deep-research-25-06-26 \\
\bottomrule
\end{tabular}
}
\caption{Correspondence between model/product names and their API identifiers.}
\label{tab:models_api_ids}
\end{table}


\subsection{Dataset Distribution}
\label{appendix_dataset_distribution}

\begin{figure*}[t]
    \centering
    \includegraphics[width=\textwidth]{figs/database_distribution1.pdf}
    \caption{Database distribution.}
    \label{fig:database_distribution}
\end{figure*}

\begin{figure*}[t]
    \centering
    \includegraphics[width=\textwidth]{figs/domain_distribution1.pdf}
    \caption{Domain distribution.}
    \label{fig:domain_distribution}
\end{figure*}
Unlike benchmarks grounded in a small set of homogeneous web sources, \textsc{SciExplore} spans over ten scientific sub-domains and integrates 16 authoritative databases with a pronounced long-tail distribution. This design forces models to navigate heterogeneous database structures and domain-specific metadata, thereby evaluating their robustness and generalization in authentic, interdisciplinary scientific workflows.

In this section, we characterize the statistical composition of \textsc{SciExplore} with respect to disciplinary scope and database coverage. 
% As illustrated in Fig.~\ref{fig:database_distribution} and Fig.~\ref{fig:domain_distribution}, these statistics collectively demonstrate that the benchmark reflects the heterogeneity and cognitive complexity of authentic scientific research workflows.

\textsc{SciExplore} encompasses more than ten scientific sub-domains spanning \textit{Life Sciences}, \textit{Chemical Sciences}, \textit{Materials Science}, \textit{Earth Sciences}, and \textit{Artificial Intelligence} (Fig.~\ref{fig:domain_distribution}). The resulting distribution reflects a deliberately balanced coverage across heterogeneous research areas, preventing over-concentration on any single discipline and enabling a more robust assessment of cross-domain generalization in scientific information-seeking tasks.

In addition, the benchmark integrates sixteen authoritative scientific databases, exhibiting a pronounced long-tail distribution (Fig.~\ref{fig:database_distribution}). While widely adopted resources such as \textit{PubChem} and \textit{Open Targets} contribute a substantial portion of instances, a significant fraction of tasks rely on highly specialized databases, including \textit{ChEMBL}, \textit{UniProt}, \textit{Materials Project}, and \textit{GWAS}. This design imposes nontrivial challenges in navigating heterogeneous database schema, access patterns, and domain-specific metadata, thereby setting \textsc{SciExplore} apart from conventional web-centric search benchmarks.

% To further quantify task difficulty, we analyze the time required for human experts to complete each task, as depicted in Fig.~\ref{fig:time_distribution}. Even for PhD-level annotators, many tasks demand substantial time due to iterative retrieval, cross-source verification, and multi-step reasoning. This distribution confirms that SciExplore predominantly consists of long-horizon, evidence-driven scientific reasoning tasks rather than simple factual lookup.


\subsection{Dataset Construction}
\label{app:datasets_construction}

The construction of \textsc{SciExplore} follows a principled, expert-driven methodology designed to faithfully reflect the cognitive demands faced by autonomous scientific research assistants. All task instances are manually crafted by domain experts to ensure realism, difficulty, and alignment with authentic scientific workflows. Rather than relying on automatic generation, our construction process emphasizes controllability, answer uniqueness, and resistance to shortcut retrieval.

\noindent{\textbf{Multi-hop Database Navigation.}}
To construct tasks requiring multi-step structured reasoning, we adopt a \emph{Reverse Trajectory Construction Strategy}. Task generation begins from a uniquely identifiable target entity and progressively replaces explicit entity mentions with composite attribute-based descriptions extracted from scientific database records. During this process, related entities are recursively masked and substituted by their own attributes, forming nested constraints that encode deep dependency chains. Consequently, the final query can only be resolved by reversing the construction trajectory through stepwise database navigation, enforcing genuine multi-hop reasoning and preventing direct lookup.

\noindent{\textbf{Ambiguous Literature Retrieval.}}
Tasks involving ambiguous literature search are constructed to simulate early-stage exploratory research under uncertainty. Experts first extract core methodological or conceptual signals from a target paper and deliberately degrade them into vague, noisy descriptions through \emph{Feature Denoising and Fuzzification}, substantially enlarging the candidate retrieval space. To ensure answer uniqueness despite this ambiguity, we introduce \emph{Validation Constraint Injection}, adding auxiliary constraints that cannot directly guide retrieval but serve as strong post-retrieval validators (e.g., publication venue or author-level properties). This design compels agents to balance broad exploration with rigorous verification.

\noindent{\textbf{Missing Reference Completion.}}
For tasks requiring citation inference, we select citation-dense paragraphs from high-impact scientific literature and systematically rewrite the surrounding context to remove surface-level textual overlap with the original references. Overly canonical or survey-style citations are excluded to avoid trivial matching. This construction emphasizes reasoning over claim–evidence relationships, ensuring that correct reference reconstruction depends on contextual understanding rather than keyword similarity.

\noindent{\textbf{Cross-Source Structured Knowledge Synthesis.}}
Tasks targeting structured synthesis are derived from expert-curated comparison tables in review papers across natural sciences and artificial intelligence. Given a predefined comparison schema, experts identify representative source papers and design tasks that require extracting heterogeneous methodological details scattered across multiple documents. Particular care is taken to include variations in terminology, implicit descriptions, and inconsistent reporting conventions. Strict schema and format constraints are enforced so that successful completion depends on accurate cross-source integration rather than partial or approximate extraction.


% \subsection{Consistency with Human Evaluation}

% \begin{table}[t]
% \centering
% \caption{Consistency with Human Evaluation for Judge Models.}
% \label{tab:judge_consistency}
% \resizebox{\columnwidth}{!}{
% \begin{tabular}{lcccc}
% \hline
% \textbf{Judge Model} & \textbf{T1 \& T2} & \textbf{T3} & \textbf{T4} & \textbf{Total} \\
% \hline
% CompassVerifier-32B      & 98.1 & 96.8 & 95.9 & 97.2 \\
% CompassJudge-32B         & 98.9 & 97.5 & 96.6 & 97.9 \\
% Qwen3-32B                & 97.7 & 96.2 & 95.1 & 96.6 \\
% Qwen2.5-Instruct-72B     & 98.8 & 97.1 & 96.0 & 97.6 \\
% GPT-OSS-120B             & \textbf{99.8} & \textbf{99.1} & \textbf{98.4} & \textbf{99.2} \\
% \hline
% \end{tabular}
% }
% \end{table}

% To verify the reliability of the automated evaluation pipeline for complex scientific tasks, we conducted a human–model consistency analysis. We randomly sampled representative model outputs and asked domain experts to manually assess their correctness against reference answers, treating expert judgments as ground truth. The same samples were then evaluated by the automated pipeline in a blind setting to measure agreement between LLM-as-a-Judge and human experts.

% As shown in Table~\ref{tab:judge_consistency}, the automated evaluation system demonstrates a high level of consistency with human judgments across all judge models. In particular, using \textit{GPT-OSS-120B}~\cite{agarwal2025gpt} as the judge yields the highest agreement score, indicating near-expert-level evaluation reliability. These results confirm that the proposed evaluation pipeline is both objective and reproducible, and can serve as a reliable substitute for costly and time-consuming manual evaluation in large-scale scientific benchmarking.

% \begin{table*}[htbp]
%     \centering
%     \caption{Statistics of Different Databases}
%     \label{tab:database_stats}
%     \resizebox{\textwidth}{!}{%
%         \begin{tabular}{lcccccccc}
%             \toprule
%             \textbf{Database} & PubChem & ChEMBL & SDBS & ORD & Open Targets & gnomAD & DrugBank & UniProt \\
%             \textbf{Num}      & 16      & 3      & 1    & 1   & 8            & 1      & 6        & 7       \\
%             \midrule
%             \textbf{Database} & mimedb & reactcome & GWAS & gnorm\_AD & dbSNP & Materials project & AFLOW & OQMD \\
%             \textbf{Num}      & 3      & 1         & 2    & 1         & 1     & 8                 & 2     & 2    \\
%             \bottomrule
%         \end{tabular}%
%     }
% \end{table*}

% \begin{table*}[htbp]
%     \centering
%     \caption{Statistics of Different Subjects}
%     \label{tab:subject_stats}
%     \resizebox{\textwidth}{!}{%
%         \begin{tabular}{lccccccc}
%             \toprule
%             \textbf{Subject} & Genomics & Proteomics & Pharmacology & Microbiology & Organic Chemistry & Analytical Chemistry & Physical Chemistry \\
%             \textbf{Num}     & 18       & 16         & 18           & 8            & 14                & 6                    & 4                  \\
%             \midrule
%             \textbf{Subject} & Crystallography & Materials Science & Geology & Oceanography & LLM & Diffusion model & \\
%             \textbf{Num}     & 5               & 10                & 2       & 3            & 4   & 3               & \\
%             \bottomrule
%         \end{tabular}%
%     }
% \end{table*}

\subsection{Evaluation Details }
\label{app:evaluation_details}

% \section{Evaluation Metrics}
% \label{sec:evaluation_metrics}

To systematically assess model performance on \textsc{SciExplore}, we design task-specific evaluation metrics tailored to the structural characteristics and reasoning requirements of each task type.
Scores are reported independently for each task, and an overall score is computed as a weighted aggregation to reflect holistic scientific information-seeking and reasoning capability.

\subsubsection{T1 and T2 Evaluation}

\noindent \textbf{T1 and T2} are evaluated using Exact Match (EM) accuracy.
A prediction is considered correct if and only if the generated answer is semantically equivalent to the ground-truth answer, allowing for minor surface-form variations (e.g., formatting or phrasing differences) that do not alter the underlying meaning.

% \begin{equation}
% F_{\text{em}} = 
% \frac{2 \cdot \text{correct}}
% {2 \cdot \text{correct} + 2 \cdot \text{incorrect} + \text{not\_attempted}},
% \end{equation}

Let $C$, $I$, and $N$ denote the numbers of correct, incorrect, and not-attempted answers, respectively. We define an Exact-Match-based score $F_{\text{em}}$ as:
\begin{equation}
F_{\text{em}} = \frac{2C}{2C + 2I + N}.
\end{equation}

This formulation penalizes incorrect and missing answers symmetrically while rewarding exact correctness.

\subsubsection{T3 Evaluation}

\noindent \textbf{T3} evaluates a model’s ability to recover supporting literature citations.
Both the ground-truth references and model-predicted references are first grouped by citation slots, and predictions are filled into a unified evaluation table for comparison.

\paragraph{Judging protocol.}
For each predicted citation entry, a judge LLM is prompted to determine whether the predicted reference is \texttt{CORRECT} or \texttt{INCORRECT}.
The judgment proceeds in two steps:
(1) whether the core cited work is correct, and
(2) whether the key bibliographic attributes (e.g., authorship, venue, or year) of the correct work are accurately identified.

\paragraph{Per-group metrics.}
For each citation group, we compute:

\begin{equation}
P = \frac{\text{correct}}{\text{correct} + \text{incorrect}},
\end{equation}

\begin{equation}
R = \frac{\text{correct}}{\text{GT\_total\_num}},
\end{equation}

\begin{equation}
F_{\text{score}} = \frac{2 \cdot P \cdot R}{P + R},
\end{equation}

where $\text{correct}$ and $\text{incorrect}$ denote the number of correctly and incorrectly predicted citations in the group, and $\text{GT\_total\_num}$ is the total number of ground-truth citations required for that group.

\paragraph{Final T3 score.}
The final citation score is obtained by averaging the $F_{\text{score}}$ across all $N$ citation groups:
\begin{equation}
F_{\text{cite}} = \frac{1}{N} \sum_{i=1}^{N} F_{\text{score}}^{(i)}.
\end{equation}

This metric jointly captures both precision (avoiding spurious references) and recall (recovering required citations), reflecting realistic citation-completion requirements.

\subsubsection{T4 Evaluation}

\noindent \textbf{T4} evaluates models on cross-source structured knowledge synthesis, where outputs are required to be presented as comparison tables.
Following prior work on structured evaluation~\cite{wong2025widesearch}, we adopt a fine-grained, multi-level evaluation protocol.

\paragraph{LLM-assisted cell judgment.}
Each predicted table cell is judged by an LLM as either \texttt{CORRECT} or \texttt{INCORRECT}.
Evaluation proceeds hierarchically: the correctness of the primary entity (row key) is assessed first, followed by the correctness of its associated attributes.

\paragraph{Item-level recall.}
To further evaluate fine-grained extraction accuracy, we compute recall over all table cells:
\begin{equation}
R_{\text{Item}} =
\frac{\text{correct}_{\text{Item}}}
{\text{GT\_total\_num}_{\text{Item}}},
\end{equation}
where $\text{correct}_{\text{Item}}$ counts correctly predicted cells (including both primary keys and attribute values), and $\text{GT\_total\_num}_{\text{Item}}$ is the total number of ground-truth cells.
Compared to PK recall, item-level recall captures both entity identification and attribute correctness.

\paragraph{Row-level recall.}
Finally, we evaluate whether complete rows are correctly reconstructed:
\begin{equation}
R_{\text{Row}} =
\frac{\text{correct}_{\text{Row}}}
{\text{GT\_total\_num}_{\text{Row}}},
\end{equation}
where a row is counted as correct only if the primary key and all associated attributes are simultaneously correct.
This strict criterion ensures that models are evaluated on their ability to produce fully consistent structured outputs.
% where $\text{correct}_{\text{cell}}$ counts all correctly predicted cells across the table.

% \paragraph{Final T4 score.}
% The final table synthesis score is defined as:
% \begin{equation}
% F_{\text{tab}} = P_{\text{tab}} \cdot R_{\text{key}}.
% \end{equation}

% This formulation emphasizes that a valid synthesis requires both correct structural coverage (row recall) and accurate attribute-level content.

\subsubsection{Overall Score}

To reflect a model’s holistic capability across scientific information-seeking and reasoning tasks of varying difficulty, we compute a weighted overall score:

\begin{equation}
\text{Score}_{\text{overall}} = \sum_{i=1}^{4} w_i \cdot \text{Score}_{T_i},
\end{equation}
where $\text{Score}_{T_i}$ denotes the task-specific score for task $T_i$ and $w_i$ is the corresponding task weight.
We adopt difficulty-aware weighting with
$w_1 = 0.2$, $w_2 = 0.2$, $w_3 = 0.3$, and $w_4 = 0.3$,
assigning higher weights to tasks that require more complex reasoning and structured synthesis.



% We employ task-specific evaluation metrics tailored to the structural characteristics of each task. T1 and T2 are evaluated using Exact Match accuracy based on semantic equivalence of the final answers. T3 is assessed with citation-level precision, measuring the proportion of correctly predicted references for each instance. For T4, we adopt a multi-granularity evaluation protocol at the row and item levels: the row-level metric measures recall over complete table rows, while the item-level metric evaluates fine-grained recall over individual cells or data items. 
% An overall score is computed as a weighted sum of the four task scores to reflect holistic model performance. 

% To systematically assess model performance on \textsc{SciExplore}, we designed task-specific metrics tailored to the structural characteristics of each task. We report independent scores for each task and calculate a weighted overall performance metric based on these results.

% \noindent \textbf{T1 and T2} adopt Exact Match (EM) accuracy as the evaluation metric. A prediction is considered correct if and only if the generated result is semantically equivalent to the ground truth.

% \noindent \textbf{T3} utilizes citation-level precision as the evaluation metric. For each sample, we calculate the ratio of correctly predicted citations to the total citations required for completion. The final T3 score is obtained by averaging these ratios across all samples.

% \noindent \textbf{T4} adopts a fine-grained evaluation protocol inspired by~\cite{wong2025widesearch} to characterize model performance in structured knowledge synthesis.Specifically, we evaluate predictions at two complementary granularities.\emph{Row-level score} treats each table row as an atomic information unit and measures the recall of correctly predicted rows with respect to the ground-truth table.\emph{Item-level score} further decomposes rows into individual cells or data items, using them as the minimum evaluation units to compute fine-grained recall.
% Together, these two metrics capture both holistic row-wise correctness and localized item-wise accuracy in generated tables.

% \noindent \textbf{Overall Score}: To reflect the model's holistic capability across different tasks, we define a weighted overall score. The final performance metric is calculated as a weighted sum of the scores from the four tasks, formulated as $\sum w_i \cdot \text{score}_i$.

\begin{table*}[t]
\centering
\caption{Performance on multi-hop T1 tasks with varying hop lengths.
All scores are percentages (\%).}
\small
\setlength{\tabcolsep}{6pt}

\begin{tabular}{l ccccc}
\toprule
\textbf{Model} &
\textbf{2-hop} & \textbf{3-hop} & \textbf{4-hop} &
\textbf{5-hop} & \textbf{6-hop} \\
\midrule

\rowcolor{gray!15}
\multicolumn{6}{l}{\textit{LLMs}} \\

DeepSeek-V3.2
& 33.33 & 20.00 & 12.50 & 12.50 & 25.00 \\

Qwen2.5-72B-Instruct
& 11.11 & 10.00 & 12.50 & 12.50 & 0.00 \\

Qwen3-30B-A3B-Thinking
& 11.11 & 20.00 & 12.50 & 12.50 & 0.00 \\

Qwen3-235B-A22B-Thinking
& 33.33 & 30.00 & 12.50 & 25.00 & 0.00 \\

GPT-5.1
& 22.22 & 20.00 & 12.50 & 12.50 & 25.00 \\

Gemini-2.5-Pro
& 33.33 & 20.00 & 12.50 & 12.50 & 25.00 \\

Gemini-3-Pro
& 33.33 & 20.00 & 12.50 & 25.00 & 25.00 \\

\midrule
\rowcolor{gray!15}
\multicolumn{6}{l}{\textit{LLMs with Search Tools}} \\

DeepSeek-V3.2 w/ search
& 22.22 & 20.00 & 12.50 & 12.50 & 25.00 \\

GPT-5.1 w/ search
& 11.11 & 30.00 & 0.00 & 12.50 & 0.00 \\

Gemini-3-Pro w/ search
& 44.44 & 40.00 & 12.50 & 25.00 & 25.00 \\

\midrule
\rowcolor{gray!15}
\multicolumn{6}{l}{\textit{Deep Research Agents}} \\

Tongyi-DeepResearch-30B-A3B
& 44.44 & 50.00 & 25.00 & 62.50 & 25.00 \\

Gemini Deep Research
& 33.33 & 20.00 & 12.50 & 25.00 & 25.00 \\

OpenAI Deep Research
& 22.22 & 40.00 & 12.50 & 12.50 & 25.00 \\

\midrule
Human
& 66.70 & 60.00 & 50.00 & 37.50 & 25.00 \\

\bottomrule
\end{tabular}
\label{tab:t1-hop-results}
\end{table*}


\subsection{Run-to-Run Variance Analysis}
\label{app:variance_analysis}

Given the moderate scale of \textsc{SciExplore}, we assess the robustness of model performance via a multi-run variance analysis. Specifically, for each evaluated model, we repeat the full benchmark evaluation four times under identical settings and compute the variance of scores across runs for each task (T1--T4) as well as the overall score. As shown in Table~\ref{tab:variance_results}, the variance is consistently low across models and tasks, with a mean overall variance of 0.78, indicating strong stability. Most models exhibit small fluctuations, suggesting that the observed performance differences are not driven by stochastic generation effects. Even for cases with relatively higher per-task variance (e.g., Gemini-2.5-Pro on T3), the overall score remains stable. These results demonstrate that our benchmark provides a reliable and reproducible evaluation signal, and that the main conclusions of the paper are robust to randomness in model outputs.

\begin{table*}[t]
\centering
\caption{Run-to-run variance across four repeated evaluations. Lower values indicate higher stability.}
\label{tab:variance_results}
\small
\begin{tabular}{lcccccc}
\toprule
\textbf{Model} & \textbf{T1} & \textbf{T2} & \textbf{T3} & \textbf{T4\_Item} & \textbf{T4\_Row} & \textbf{Overall} \\
\midrule
DeepSeek-V3.2 & 2.10 & 2.40 & 2.25 & 1.05 & 0.67 & 1.41 \\
Qwen2.5-72B-Instruct & 0.34 & 0.43 & 0.52 & 0.52 & 0.24 & 0.24 \\
Qwen3-30B-A3B-Thinking & 1.02 & 0.18 & 0.55 & 0.06 & 0.14 & 0.64 \\
Qwen3-235B-A22B-Thinking & 1.17 & 3.24 & 3.19 & 1.89 & 0.44 & 0.64 \\
GPT-5.1 & 2.75 & 2.43 & 2.90 & 1.67 & 0.39 & 0.21 \\
Gemini-2.5-Pro & 2.30 & 0.00 & 4.02 & 0.82 & 0.12 & 1.32 \\
Gemini-3-Pro & 1.65 & 3.24 & 0.85 & 0.73 & 0.29 & 0.28 \\
DeepSeek-V3.2 w/ search & 0.36 & 0.20 & 1.16 & 0.18 & 0.39 & 0.45 \\
GPT-5.1 w/ search & 2.85 & 2.60 & 2.95 & 4.86 & 1.61 & 2.91 \\
Gemini-3-Pro w/ search & 0.98 & 0.22 & 0.66 & 0.29 & 0.12 & 0.76 \\
Tongyi-DeepResearch-30B-A3B & 0.53 & 0.29 & 0.67 & 0.51 & 0.31 & 0.28 \\
Gemini Deep Research & 0.59 & 0.29 & 0.18 & 0.09 & 0.35 & 0.36 \\
OpenAI Deep Research & 0.92 & 0.93 & 0.80 & 0.13 & 0.07 & 0.68 \\
\midrule
\textbf{Mean} & 1.35 & 1.27 & 1.59 & 0.98 & 0.40 & 0.78 \\
\bottomrule
\end{tabular}
\end{table*}


\subsection{Error Analysis}
\label{app:error_analysis}

This section provides detailed examples for each of the failure modes identified in Section 5. We utilize the trajectories of representative state-of-the-art agents (e.g., Gemini-DeepResearch or GPT-4o) to illustrate these distinct cognitive breakdowns. Each case includes the user’s task, the agent’s actions (queries and retrieval), and a diagnostic analysis of the specific error.

\begin{figure*}[t]
    \centering
    \includegraphics[width=\linewidth]{figs/case_study_Premature.pdf}
    \caption{Premature Abandonment due to Search Inertia. The agent terminates the search after failing to find a direct semantic match, ignoring alternative search strategies.
    % This example demonstrates that if the model fails to retrieve information relevant to the answer during the initial search phase, it tends to abandon further retrieval immediately, rather than attempting to adjust its search strategy or employ more diverse queries to continue.
    }
    \label{fig:e1}
\end{figure*}

\noindent\textbf{Case 1: Premature Abandonment due to Search Inertia} Figure \ref{fig:e1} illustrates a case of premature abandonment when facing strict retrieval constraints. The agent was tasked with identifying specific mutations in the RNLS gene associated with "cerebral amyloid deposition" with a highly specific P-value threshold ($< 8 \times 10^{-7}$) and requesting allele frequencies from the gnomAD database. The agent initiated a reasonable keyword search ("RNLS gene cerebral amyloid deposition..."). However, upon finding no direct semantic match in the first page of search results, the agent immediately concluded that no such evidence exists and terminated the process. This demonstrates a lack of strategic resilience; instead of broadening the search scope (e.g., searching for the phenotype in GWAS catalogs first, then cross-referencing the gene) or attempting synonymous queries, the agent exhibited "search inertia," treating the absence of immediate results as proof of non-existence.

\begin{figure*}[t]
    \centering
    \includegraphics[width=\linewidth]{figs/case_study_hallucinated_response.pdf} 
    \caption{Hallucination in Constraint Verification. The agent fabricates details (reference count and number of spectra) to force alignment with the user's complex constraints.}
    \label{fig:e2}
\end{figure*}

\noindent\textbf{Case 2: Hallucination in Constraint Verification} Figure \ref{fig:e2} provides an example of model hallucination triggered by complex constraint verification. The task required identifying a specific paper published after 2015, containing over 50 references, and featuring five spectra in Figure 6. The agent successfully retrieved a candidate paper ("A Plasmonic Coupling Substrate...") that matched the topic and publication year. However, it failed to rigorously verify the fine-grained visual and structural constraints. Instead of rejecting the paper or performing a specific "Ctrl+F" style verification for the reference count, the agent hallucinated that the paper contained "more than 50 references" and "five spectra" to align with the user's prompt. This error highlights a critical reliability gap where the model prioritizes user-compliance over factual groundedness, fabricating details to force a retrieved document to fit the query criteria.

\begin{figure*}[t]
    \centering
    \includegraphics[width=\linewidth]{figs/case_study_long_context.pdf}
    \caption{Information Loss in Long-Context Extraction. The agent relies on the abstract summary ("20 subjects") rather than the ground truth buried in the results ("18 subjects").}
    \label{fig:e3}
\end{figure*}

\noindent\textbf{Case 3: Information Loss in Long-Context Extraction} Figure \ref{fig:e3} illustrates a failure in precise evidence extraction within long-context scenarios ("Needle-in-a-Haystack" failure). The task involved summarizing tDCS studies, specifically requiring the number of samples/subjects. The agent correctly located the relevant study (Khan et al., 2020) and ingested the full text. However, it extracted the sample size as "20" based on the Abstract ("Twenty subjects were randomly assigned..."). It failed to attend to the conflicting, ground-truth information buried deeper in the Results or Subjects section, which clarified that "only 18 subjects... were able to finish." This discrepancy indicates that despite having large context windows, current models exhibit a attention bias toward high-level summaries (abstracts) and struggle to resolve conflicting information dispersed across long scientific texts.

\begin{figure*}[t]
    \centering
    \includegraphics[width=\linewidth]{figs/case_study_output_structure.pdf}
    \caption{Structural Instruction Non-Compliance. The agent omits required columns ("Key Chemical Properties" and "Examples"), failing to adhere to the strict table schema provided.}
    \label{fig:e4}
\end{figure*}

\noindent\textbf{Case 4: Structural Instruction Non-Compliance} Figure \ref{fig:e4} presents a failure in following strict structural constraints during synthesis. The task explicitly required the generation of a Markdown table with five specific headers: "Polymer Type," "Key Chemical Properties," "Drug Release Mechanism," "Applications," and "Examples." While the agent successfully retrieved relevant content regarding polymeric drug delivery systems, the final output completely omitted the "Key Chemical Properties" and "Examples" columns, collapsing the information into a simplified three-column table. This exemplifies a schema alignment failure, where the model's internal training bias toward generating generic summaries overrides the explicit formatting constraints provided in the prompt. In automated scientific workflows where downstream parsers rely on fixed schemas, such structural instability renders the agent unusable.

\noindent\textbf{Case 5: Premature Hypothesis Commitment under Search.
} 
Figure~\ref{fig:e5} shows a failure case where search augmentation degrades performance by amplifying a hypothesis-first reasoning bias. Although the task explicitly specified a hard constraint on the target protein length (604 aa), the model prematurely inferred the pain-associated sodium channel Nav1.8 (SCN10A) based on superficial semantic cues. When search results later indicated that the canonical SCN10A protein is approximately 1956 aa, the model failed to revise its hypothesis and instead fabricated a non-existent “604 aa SCN10A construct” to satisfy the constraint. This behavior reflects post-hoc rationalization rather than evidence-driven hypothesis updating, where retrieved information is selectively distorted to justify an internally committed conclusion. Such failures highlight a key limitation of current search-augmented LLMs: external retrieval does not guarantee improved correctness when constraint enforcement and hypothesis revision mechanisms are absent.


\begin{figure*}[t]
    \centering
    \includegraphics[width=\linewidth]{figs/case_study_gpt_search_or_no.pdf}
    \caption{Comparison between search-enabled and non-search responses on a constraint-sensitive biomedical query. While the non-search model correctly identifies Prostaglandin G/H synthase 2 by adhering to the explicit sequence-length constraint (604 aa), the search-enabled model prematurely commits to a pain-related hypothesis (Nav1.8 / SCN10A) and fails to revise it after encountering contradictory evidence, resulting in an incorrect answer.}
    \label{fig:e5}
\end{figure*}

\begin{figure*}[t]
    \centering
    \includegraphics[width=\linewidth]{figs/case_study_thinking_process.pdf} 
    \caption{Illustration of hypothesis-driven reasoning during model inference.
The model first proposes candidate interpretations for Compound B based on partial chemical cues, and subsequently validates these hypotheses against reaction mechanisms and known enzyme functions. The final answer is obtained after this verification stage, exemplifying a hypothesis-driven (guess-then-check) problem-solving process.}
    \label{fig:e6}
\end{figure*}

\subsection{Representative QA Examples of \textsc{SciExplore}}
\label{app:qa_examples}

% In this section, we present representative question--answer examples for each task type (T1--T4) in \textsc{SciExplore}.
In this section, we present representative question--answer examples for each task type (T1--T4) in \textsc{SciExplore}, illustrating the diverse reasoning and information-seeking behaviors required by the benchmark. These examples highlight how tasks progressively increase in complexity, ranging from structured database navigation and ambiguous literature retrieval to missing citation completion and cross-source structured knowledge synthesis. By examining concrete model inputs and outputs, we aim to provide an intuitive understanding of the challenges posed by each task type and the typical failure modes observed in current systems. Representative examples for T1, T2, T3, and T4 are shown in Figures~\ref{fig:example-t1},~\ref{fig:example-t2},~\ref{fig:example-t3},~\ref{fig:example-t4}, respectively.

\subsection{Prompt Templates}
We present the prompt templates used for each component in the \textsc{SciExplore} evaluation pipeline at the end of Appendix.

\newcounter{promptbox}
\renewcommand{\thepromptbox}{P\arabic{promptbox}}

\clearpage
\begin{figure*}[t]
\centering
\begin{examplebox}
\textbf{Question.}  
The molecular formula of Compound A is C$_9$H$_{11}$N$_2$O$_{14}$P$_3^{4-}$, and its Topological Polar Surface Area (TPSA) is 250~\AA$^2$.  
Among the chemical substances involved in the pathway associated with \textit{Aquilegia coerulea} (Colorado blue columbine) is Compound B.  
The molecular formula of Compound B is C$_9$H$_{12}$N$_3$O$_7$P$^{2-}$, and it is one of the products of the following biochemical reaction:  
2'-deoxycytidine reacts with a ribonucleoside 5'-triphosphate to yield H$^+$, Compound B, and a ribonucleoside 5'-diphosphate.  
The enzyme catalyzing this reaction in mouse is Protein C.  
What is the sequence length of Protein C?  
[The prediction error range is $\pm 0$]

\vspace{0.5em}
\textbf{Answer.}  
260
\end{examplebox}
\caption{\textbf{T1 Example Scientific Database Navigation}.}
\label{fig:example-t1}
\end{figure*}

\begin{figure*}[t]
\centering
\begin{examplebox}
\textbf{Question.}  
A research paper obtained various mixed crystals by adjusting the feed ratio.  
In the experiments of the paper, the surface growth method was used.  
The paper contains six figures in total, was published after 2018, and has at least one author with the surname Wang.  
What is the title of the paper?

\vspace{0.5em}
\textbf{Answer.}  
A solid-solution approach for controllable photomechanical crystalline materials
\end{examplebox}
\caption{\textbf{T2 Example Ambiguous Literature Retrieval}.}
\label{fig:example-t2}
\end{figure*}

\begin{figure*}[t]
\centering
\begin{examplebox}
\textbf{Question.}  
I would like to write an article on plastic pollution.  
The introduction section is already complete, but the sources supporting some key conclusions are still missing—please help me add them.

Plastic chemicals may also impede the transition to a circular economy, including technological solutions to plastic pollution[ref.1].  
For instance, increasing the reuse of plastic products can increase the release of chemicals[ref.2], and uncontrolled recycling can further perpetuate the spread of hazardous chemicals in products with high exposure potential, such as food packaging and toys.  
Chemicals in plastics can also hinder sorting and recycling and thereby impede the production of high-quality secondary materials[ref.3].  

Previous efforts to compile an overview of the chemical diversity of plastics[ref.4] have jointly identified more than 13,000 known plastic chemicals.

\textbf{Problem description:}  
All in-text citations appear as ``[ref.NUMBER]''.  
Please supply the complete bibliographic information for every citation, adhering to the specifications below.

\vspace{0.5em}
\textbf{Answer.}  

[ref.1: Wang, Z. \& Praetorius, A. Integrating a chemicals perspective into the global plastic treaty. Environ. 2022.]  

[ref.2: Mannoni, V. et al. Migration of formaldehyde and melamine from melaware and other amino resin tableware in real life service. \textit{Food Addit. Contam.} 2017.]  

[ref.3: Wiesinger, H., Wang, Z. \& Hellweg, S. Deep dive into plastic monomers, additives, and processing aids. 2021.]  

[ref.4: Aurisano, N., Weber, R. \& Fantke, P. Enabling a circular economy for chemicals in plastics. 2021.]
\end{examplebox}
\caption{\textbf{T3 Example Missing Citation Completion}.}
\label{fig:example-t3}
\end{figure*}

\begin{figure*}[t]
\centering
\begin{examplebox}
\textbf{Question.}

Find representative LLMs applied in bioinformatics in 2024, and list a table showing their Model name, Venue name, Model type, and Task type.

Here, the Model type includes three categories: Encoder-only, Decoder-only, Encoder-Decoder; if none of these apply, use - to indicate.
The Task type includes four categories: DNA, RNA, protein, scRNA.

Please organize the table in Markdown format.

\vspace{0.5em}
\textbf{Answer.}

\begingroup
\footnotesize
\setlength{\tabcolsep}{3pt}
\begin{verbatim}
| Model                  | Venue                           | Type            | Task    |
|------------------------|---------------------------------|-----------------|---------|
| RNA-MSM                | Nucleic Acids Research          | Encoder-only    | RNA     |
| GROVER                 | Nature Machine Intelligence     | Encoder-only    | DNA     |
| megaDNA                | Nature Communications           | Decoder-only    | DNA     |
| Nucleotide Transformer | Nature Methods                  | Encoder-only    | RNA     |
| RNAErnie               | Nature Machine Intelligence     | Encoder-only    | RNA     |
| RhoFold+               | Nature Methods                  | Encoder-only    | RNA     |
| GPTCelltype            | Nature Methods                  | Decoder-only    | scRNA   |
| Evo                    | Science                         | Decoder-only    | DNA     |
| scFoundation           | Nature Methods                  | Encoder-Decoder | scRNA   |
| scGPT                  | Nature Methods                  | Encoder-Decoder | scRNA   |
| AlphaFold3             | Nature                          | -               | Protein |
| ESM-DBP                | Nature Communications           | Encoder-only    | Protein |
| RoseTTAFold All-Atom   | Science                         | Encoder-Decoder | Protein |
| ProstT5                | NAR Genomics and Bioinformatics | Encoder-Decoder | Protein |
| EpiGcPT                | NAR Genomics and Bioinformatics | Encoder-only    | DNA     |
| RiNALMo                | arXiv                           | Encoder-only    | RNA     |
| EMBED                  | Bioinformatics Advances         | Encoder-Decoder | DNA     |
\end{verbatim}
\endgroup
\end{examplebox}
\caption{\textbf{T4 Example Cross-Source Structured Knowledge Synthesis
}.}
\label{fig:example-t4}
\end{figure*}


\onecolumn

\refstepcounter{promptbox}

\begin{tcolorbox}[
    breakable,             
    % enhanced,   
    % enhanced jigsaw,
    rounded corners,
    width=\textwidth,
    title={T1 and T2 JUDGE PROMPT},
    fonttitle=\normalsize\bfseries,
    coltitle=white,
    toptitle=2mm, bottomtitle=2mm, left=3mm, right=3mm, top=3mm, bottom=3mm,
]

\subsection*{Task Description}

As a grading expert, your task is to determine whether the candidate's final answer matches the provided standard answer. You must distinguish between exact matches, acceptable variations, and incorrect responses based on the guidelines below.

\subsection*{Evaluation Protocol}

Please follow these guidelines precisely:

\begin{enumerate}
    \item \textbf{Reference Standard:}
    \begin{itemize}
        \item The standard answer is definitive but allows for a certain margin of error.
        \item Do not regenerate answers; only compare with the given standard.
    \end{itemize}

    \item \textbf{Comparison Method:}
    \begin{itemize}
        \item Carefully analyze the question's requirements and the standard answer's structure:
        \begin{itemize}
            \item Determine whether the question expects \textbf{exact matching} or allows \textbf{partial matching}.
            \item Base this determination on the question's phrasing and the nature of the standard answer.
        \end{itemize}
        \item Compare \textbf{ONLY} the candidate's final answer (ignore reasoning/explanation errors).
        \item Disregard formatting or presentation style differences.
        \item Use discretion for unit omission (avoid automatic rejection).
        \item \textbf{Math:} Calculate step-by-step to check if formulas are equivalent.
        \item \textbf{Multiple Choice:} Compare only the final choice and option content.
    \end{itemize}

    \item \textbf{Multi-part Answers:}
    \begin{itemize}
        \item For questions requiring multiple responses (e.g., multi-select): All parts must match the standard answer exactly.
        \item Compare each sub-answer step by step. Partial matches are considered incorrect.
    \end{itemize}

    \item \textbf{Validity Check:}
    Reject answers that meet the following criteria (specify type in judgment):
    \begin{itemize}
        \item \textbf{Incomplete:} Cut off mid-sentence or lacks a complete response $\rightarrow$ Label as \textbf{INCOMPLETE}.
        \item \textbf{Repetitive:} Loops words or phrases $\rightarrow$ Label as \textbf{REPETITIVE}.
        \item \textbf{Refusal:} Explicit refusals (e.g., "I cannot answer...") $\rightarrow$ Label as \textbf{REFUSAL}.
    \end{itemize}
\end{enumerate}

\subsection*{Grading Scale}

\begin{itemize}
    \item $\boxed{A}$ \textbf{- CORRECT:}
    \begin{itemize}
        \item Answer matches standard exactly (including equivalent expressions).
        \item Within the allowable error range specified by the question.
        \item Semantically equivalent responses.
    \end{itemize}

    \item $\boxed{B}$ \textbf{- INCORRECT:}
    \begin{itemize}
        \item Outside the allowable error range specified by the question.
    \end{itemize}

    \item $\boxed{C}$ \textbf{- INCOMPLETE/REPETITIVE/REFUSAL:}
    \begin{itemize}
        \item Fails validity criteria (must specify subtype).
    \end{itemize}
\end{itemize}

\section*{Output Format}

Please strictly use the following format for your response:

\begin{lstlisting}[breaklines=true, basicstyle=\ttfamily]
Analysis step by step: [
Thoroughly evaluate the candidate's answer including:
(1) First check if the answer is INCOMPLETE, REPETITIVE, or a REFUSAL - if so, immediately classify as \boxed{C} with the corresponding type.
(2) Analyze the question's core requirements and the standard answer's structure (strict requirements vs. tolerant allowances).
(3) Perform a detailed comparison between the candidate's final answer and the standard answer (content equivalence, numerical precision, expression formats).
]
Final Judgment: \boxed{A/B/C} - <CORRECT/INCORRECT/INCOMPLETE/REPETITIVE/REFUSAL>
\end{lstlisting}

\subsection*{Task Input}

\texttt{<Original Question Begin>}\\
\texttt{\{question\}}\\
\texttt{<Original Question End>}

\vspace{0.2cm}

\texttt{<Standard Answer Begin>}\\
\texttt{\{gold\_answer\}}\\
\texttt{<Standard Answer End>}

\vspace{0.2cm}

\texttt{<Candidate's Answer Begin>}\\
\texttt{\{llm\_response\}}\\
\texttt{<Candidate's Answer End>}

\end{tcolorbox}

\begin{tcolorbox}[
    breakable,              % 允许跨页
    enhanced,
    rounded corners,
    width=\textwidth,
    title={T4 JUDGE PROMPT},
    fonttitle=\normalsize\bfseries,
    coltitle=white,
    toptitle=2mm, bottomtitle=2mm, left=3mm, right=3mm, top=3mm, bottom=3mm,
]

\subsection*{Role}

You are a strict evaluator for \textbf{single-row item-level recall} against a predicted table.

\subsection*{Inputs}

You will receive:

\begin{itemize}
    \item \textbf{Gold Row Table (ANSWER)}: Contains \textbf{a header + exactly one data row}.\\
    \texttt{\{ANSWER\_TABLE\}}
    
    \item \textbf{Predicted Table (PREDICTION)}: Contains a header + multiple rows (or may be empty).\\
    \texttt{\{PREDICTION\_TABLE\}}
\end{itemize}

\subsection*{Objective}

Determine whether PREDICTION contains the gold row (by primary key), then compute \textbf{item-level recall} for that single gold row.

\begin{itemize}
    \item The \textbf{first column} of ANSWER is the \textbf{primary key}.
    \item If the gold primary key is \textbf{not present} in any prediction row (after normalization), then \textbf{recall = 0.0}.
    \item If the primary key \textbf{is present}, compare the remaining cells (items) and compute:
\end{itemize}

Let:
\begin{itemize}
    \item \texttt{items\_total} = number of gold cells in the gold row \textbf{excluding} the primary key column, excluding empty/NA.
    \item \texttt{items\_recalled} = number of those cells that are matched by the predicted row with the same primary key (after column mapping + normalization).
\end{itemize}

Then:
\[ \texttt{recall} = \frac{\texttt{items\_recalled}}{\texttt{items\_total}} \]

Return \texttt{recall} as a float (e.g., 0.2) plus minimal supporting details.

\subsection*{Normalization Rules (apply BEFORE matching)}

\begin{enumerate}
    \item Ignore case, extra whitespace, and trivial punctuation.
    \item \textbf{Numeric tolerance}:
    \begin{itemize}
        \item Match if absolute diff $\le 10^{-6}$ OR relative diff $\le 10^{-3}$ (when $|\text{value}| > 10^{-6}$).
        \item Treat \texttt{0.25} as equal to \texttt{25\%} if clearly percentage-formatted.
    \end{itemize}
    \item \textbf{Units}: convert when unambiguous (e.g., g $\leftrightarrow$ mg). If units conflict, treat as mismatch.
    \item \textbf{Aliases}: allow common aliases if clearly the same entity.
    \item \textbf{List-in-cell}:
    \begin{itemize}
        \item If a cell contains a list separated by commas/semicolons, treat as a set.
        \item Count as recalled only if predicted set contains all gold elements (subset match).
    \end{itemize}
    \item \textbf{Be conservative}: if uncertain, mark mismatch.
\end{enumerate}

\subsection*{Column Mapping}

\begin{itemize}
    \item Use ANSWER headers as the canonical schema.
    \item Map each gold column to the best matching prediction column by:
    \begin{enumerate}
        \item exact header match after normalization, else
        \item semantic similarity of header names.
    \end{enumerate}
    \item If a gold column cannot be mapped, that item is \textbf{not recalled}.
\end{itemize}

\subsection*{Matching Procedure (must follow)}

\begin{enumerate}
    \item Parse ANSWER: extract headers and the single gold row.
    \item Identify $\texttt{pk\_col} = \text{first header}$, $\texttt{gold\_pk} = \text{normalized}(\text{value in pk\_col})$.
    \item Parse PREDICTION rows and build an index from normalized pk values to candidate rows.
    \item If \texttt{gold\_pk} not found in prediction index:
    \begin{itemize}
        \item Set \texttt{items\_total} as defined above.
        \item Set $\texttt{items\_recalled} = 0$.
        \item Set $\texttt{recall} = 0.0$.
    \end{itemize}
    \item If \texttt{gold\_pk} found:
    \begin{itemize}
        \item If multiple predicted rows share the same pk, choose the row that maximizes \texttt{items\_recalled}.
        \item For each eligible gold item cell:
        \begin{itemize}
            \item Compare gold cell vs predicted mapped cell after normalization.
            \item Count match as recalled; otherwise mismatch.
        \end{itemize}
    \end{itemize}
\end{enumerate}

\section*{Output Format (STRICT JSON ONLY)}

Return ONLY one JSON object:

\begin{lstlisting}[breaklines=true, basicstyle=\ttfamily, frame=none]
{
  "recall": float,
  "items_recalled": int,
  "items_total": int,
  "primary_key_column": "string",
  "gold_primary_key": "string (normalized)",
  "matched_prediction_row_index": int | null,
  "column_mapping": { "gold_col": "pred_col_or_null", ... },
  "mismatches": [
    {
      "column": "gold column name",
      "gold_value": "normalized gold cell",
      "pred_value": "normalized predicted cell or null",
      "reason": "short reason"
    }
  ]
}
\end{lstlisting}

\subsection*{Important Constraints}

\begin{itemize}
    \item Do not output any text outside the JSON.
    \item Do not hallucinate missing values.
    \item If $\texttt{items\_total} == 0$, set $\texttt{recall} = 0.0$.
\end{itemize}

\end{tcolorbox}